% Шаблон статьи article-forge. Сборка: python -m forge build <этот каталог>.
\documentclass[10pt,a4paper,twocolumn]{article}
\usepackage{artforge}            % [english] — английский документ

\input{numbers.tex}              % \N<Key> — числа из figures.py

\ArticleTopLine{Проект\ \textbullet\ Технический отчёт TR-0000-00}{Месяц 2026}
\ArticleRunningHead{Проект · Технический отчёт TR-0000-00}{Краткое название}
\ArticleTitle{Заголовок статьи:\\ что сделано и что это даёт}
\ArticleAuthors{Автор\textsuperscript{1}\ \ \textbullet\ \ Соавтор\textsuperscript{2}}
\ArticleAffiliations{\textsuperscript{1}Организация\quad\textsuperscript{2}Организация}
\ArticleAbstract{%
Одна-две фразы о задаче. Что сделано. Главный измеренный результат: медианное время снизилось
с \NBaseMs\,мс до \NBestMs\,мс, в \NSpeedup~раза (\NRuns{} прогонов).}
\ArticleKeywords{ключевое слово; ещё одно.}
\ArticleHighlights{%           % необязательно: без него аннотация занимает всю ширину
\NSpeedup\texttimes & ускорение по медиане\\
\NRuns{} & прогонов в эксперименте\\
}

\begin{document}
\MakeArticleTitle

\section{Введение}
Постановка задачи и вклад работы. Идентификаторы из кода набираются через \code{module.function_name}~—
они переносятся в любом месте.

\begin{keybox}{Главное}
Вывод в одну фразу, который читатель должен унести.
\end{keybox}

\section{Методика}
Как проводились измерения, на каком окружении, сколько повторов.

\section{Результаты}
Рис.~\ref{fig:speed} показывает медианное время для трёх вариантов.

\begin{figure}[t]
\centering\includegraphics[width=\linewidth]{figures/fig_speed.pdf}
\caption{Медианное время обработки по вариантам; \NRuns{} прогонов.}
\label{fig:speed}
\end{figure}

\begin{table}[t]
\caption{Пример таблицы.}
\centering\tabfont
\begin{tabular}{@{}lr@{}}
\toprule
Вариант & Медиана, мс\\
\midrule
базовый & \NBaseMs\\
лучший & \NBestMs\\
\bottomrule
\end{tabular}
\end{table}

\section{Обсуждение}
Ограничения и что дальше. Решения, не подтверждённые источником, помечаются \decision.

\balance
\begin{thebibliography}{9}
\bibitem{ex} Автор. Название. Издание, год. \url{https://example.org/}
\end{thebibliography}

\end{document}
